\chapter{Wykład X. Wstęp do Teorii Pierścieni.}

\section{Definicja Pierścienia}

\begin{remark}
	Mówimy, że zbiór $A$ wraz z działaniem $\cdot $ jest grupą (oznaczamy $\left( A, \cdot  \right) $ ), jeżeli spełnione są następujące aksjomaty:
\begin{enumerate}[label=A\arabic*.]
	\item Łączność: $\left(\forall a,b,c \right)\left( a\cdot \left( b\cdot c \right) = \left( a\cdot b \right) \cdot c \right)  $.
	\item Istnienie elementu neutralnego: $\left(\exists e in A \right) \left(\forall a \in A \right)\left( a\cdot e = e\cdot a = a \right) $.
	\item Istnienie elementu odwrotnego dla każdego elementu grupy: $\left(\forall a \in A \right)\left(\exists b \in A \right)\left( a\cdot b= b\cdot a = e \right)$.
\end{enumerate}
\end{remark} 
Zaczynamy od ustalenia zbioru $R$ z dwoma działaniami (wewnętrznymi): dodawaniem ($+$) i mnożeniem ($\cdot$).
\begin{align*}
    +: R \times R &\to R \\
    \cdot: R \times R &\to R
\end{align*}
\begin{remark}
	Przyjmujemy formalne ustalenie kolejności wykonywania działań: $x \cdot y + z := (x \cdot y) + z$. Jest to niezbędne dla uniknięcia niejednoznaczności w dowodach formalnych
\end{remark} 

\begin{definition}[Aksjomaty Pierścienia]
    Strukturę $(R, +, \cdot)$ nazywamy \textbf{pierścieniem}, jeśli spełnione są następujące warunki (oznaczmy je A1-A3):
    \begin{enumerate}[label=A\arabic*.]
        \item $(R, +)$ jest \textbf{grupą abelową} (przemienną).
        \begin{itemize}
            \item Element neutralny dodawania oznaczamy przez $\mathbf{0}$ (zero pierścienia).
            \item Element przeciwny do $a$ oznaczamy przez $-a$.
        \end{itemize}
        \item \textbf{Rozdzielność} mnożenia względem dodawania:
        \[ (\forall a,b,c \in R) \left( a \cdot (b+c) = a \cdot b + a \cdot c \quad \land \quad (b+c) \cdot a = b \cdot a + c \cdot a \right) \]
        \item \textbf{Łączność} mnożenia:
        \[ (\forall a,b,c \in R) \left( a \cdot (b \cdot c) = (a \cdot b) \cdot c \right) \]
    \end{enumerate}
    
    Dodatkowo wyróżniamy specjalne typy pierścieni, jeśli spełniają kolejne aksjomaty:
    \begin{enumerate}[label=A\arabic*., start=4]
        \item \textbf{Istnienie jedynki}:
        \[ (\exists \mathbf{1} \in R)(\forall a \in R) (a \cdot \mathbf{1} = \mathbf{1} \cdot a = a) \]
        Jeśli zachodzi A1-A4, mówimy o \textbf{pierścieniu z jednością} (lub z jedynką).
        
	\textit{Uwaga: Z aksjomatów łączności (A3) i istnienia jedynki (A4) wynika, że jedynka jest wyznaczona jednoznacznie}
        \item \textbf{Przemienność mnożenia}:
        \[ (\forall a,b \in R) (a \cdot b = b \cdot a) \]
        Jeśli zachodzi A1-A3 oraz A5, mówimy o \textbf{pierścieniu przemiennym}.
    \end{enumerate}
\end{definition}

\noindent \textbf{Podsumowanie nazewnictwa:}
\begin{itemize}
    \item \textbf{Pierścień:} Spełnia A1, A2, A3.
    \item \textbf{Pierścień z jednością:} Spełnia A1, A2, A3, A4.
    \item \textbf{Pierścień przemienny:} Spełnia A1, A2, A3, A5.
    \item \textbf{Pierścień przemienny z jednością:} Spełnia wszystkie warunki A1-A5. (To jest nasz główny obiekt zainteresowań w dalszej części kursu).
\end{itemize}

\begin{remark}[Konwencje zapisu]
    W dalszej części, pisząc "pierścień $R$", będziemy stosować skróty:
    \begin{itemize}
        \item $ab := a \cdot b$ (opuszczanie kropki),
        \item $a^n := \underbrace{a \cdot \ldots \cdot a}_{n \text{ razy}}$ dla $n \in \mathbb{N}_{>0}$,
        \item $a - b := a + (-b)$.
    \end{itemize}
\end{remark}

\subsection*{Podpierścienie i Przykłady}

\begin{definition}[Podpierścień]
    Niech $(R, +, \cdot)$ będzie pierścieniem. Podzbiór $S \subseteq R$ nazywamy \textbf{podpierścieniem} (ozn. $S \underset{\text{podp}}{\subseteq} R$), gdy:
    \begin{enumerate}
        \item $S$ jest podgrupą grupy addytywnej $(R, +)$ (czyli $S \le (R,+)$),
        \item $S$ jest zamknięty na mnożenie: $(\forall x,y \in S)(xy \in S)$.
    \end{enumerate}
    Wtedy $(S, +, \cdot)$ sam w sobie jest pierścieniem (dziedziczy własności z $R$).
\end{definition}
\begin{remark}
Oczywiście $(S, +\restriction_{\scriptscriptstyle S\times S}, \cdot\restriction_{\scriptscriptstyle S\times S})$ jest pierścieniem, gdy $S$ jest podpierścieniem $R$.
\end{remark} 

\subsubsection*{Katalog przykładów}

\begin{enumerate}[label=\arabic*)]
    \item $(\mathbb{Z}, +, \cdot)$: Wzorcowy przykład pierścienia przemiennego z jedynką $\mathbf{1}=1$.
    
    \item $(\mathbb{N}, +, \cdot)$: \textbf{NIE} jest pierścieniem (nie jest grupą ze względu na dodawanie – brak elementów przeciwnych).
    
    \item $(2\mathbb{Z}, +, \cdot)$: Zbiór liczb parzystych. Jest to pierścień przemienny, ale \textbf{bez jedynki} (bo $1 \notin 2\mathbb{Z}$). Jest to podpierścień $\mathbb{Z}$.
    
    \item $(\mathbb{Z}_n, +_n, \cdot_n)$: Pierścień reszt modulo $n$. Jest przemienny z jedynką.
    
    \item Liczby wymierne $\mathbb{Q}$, rzeczywiste $\mathbb{R}$, zespolone $\mathbb{C}$: Wszystkie są pierścieniami przemiennymi z jedynką. Mamy ciąg inkluzji podpierścieni:
    \[ 2\mathbb{Z} \underset{\text{podp}}{\subseteq} \mathbb{Z} \underset{\text{podp}}{\subseteq} \mathbb{Q} \underset{\text{podp}}{\subseteq} \mathbb{R} \underset{\text{podp}}{\subseteq} \mathbb{C} \]
    
    \item $\left( M_{n} \left( \mathbb{R}  \right), +, \cdot   \right) $ jest pierścieniem z $\mathbf{1}$. Ponadto jest przemienny $\iff n = 1$.

    \item \textbf{Pierścienie macierzy} $M_n(R)$:
    Niech $R$ będzie pierścieniem przemiennym z $\mathbf{1}$. Zbiór macierzy kwadratowych $n \times n$ o wyrazach z $R$ tworzy pierścień wraz z działaniami z $R$. Dokładniej to tak samo jak w $M_{n}\left( \mathbb{R}  \right) $ definiujemy dodawanie i mnożenie macierzy: używamy $+$ oraz  $*$ z $R$. 
    \begin{itemize}
        \item Jedynką jest macierz diagonalna $\mathcal{I} = \text{diag}(\mathbf{1}_R, \dots, \mathbf{1}_R)$:
	\[
		\mathbf{1} = \mathcal{I} = \begin{pmatrix} \mathbf{1}_{R} & \ldots & 0 \\ \vdots & \ddots & \vdots \\ 0 & \ldots & \mathbf{1}_{R} \end{pmatrix} 
	.\] 
        \item \textbf{Uwaga:} Jeśli $n > 1$, ten pierścień jest \textbf{nieprzemienny} (nawet jeśli $R$ jest przemienny).
    \end{itemize}
    \textbf{Przykład:} $M_3\left( \mathbb{Z} _{3} \right) $ jest pierścieniem. 
    
    \item \textbf{Pierścienie macierzy} $M_n(R)$:
Niech $R$ będzie dowolnym pierścieniem przemiennym z $\mathbf{1}$. Zdefiniowanie dodawania i mnożenia macierzy pozwala stworzyć pierścień $M_n(R)$, np. $M_3(\mathbb{Z}_3)$.
    \item \textbf{Pierścienie funkcyjne} $\text{Fun}(X, R)$:
    Niech $X$ będzie dowolnym zbiorem, a $R$ pierścieniem. Definiujemy $\text{Fun}\left( X, R \right) $ jako zbiór wszystkich funkcji $f: X \to R$ (oznaczany też $R^X$) z działaniami "po punktach":
    \begin{align*}
        (f+g)(x) &:= f(x) +_R g(x) \\
        (f \cdot g)(x) &:= f(x) \cdot_R g(x)
    \end{align*}
    Tak zdefiniowane $\text{Fun}\left( X, R \right) $ jest pierścieniem.

    Ponadto
    \begin{itemize}
        \item Jeśli $R$ ma jedynkę, to $\text{Fun}(X,R)$ też ma (funkcja stała równa $\mathbf{1}_R$).
	\item Jeśli $R$ jest przemienne, to $\text{Fun}\left( X, R \right) $ też jest przemienny.
    \end{itemize}
    \item Definiujemy podpierścienie $\text{Fun}\left( R, R \right)$:
		\begin{itemize}
			\item $C\left( R \right) $ : pierścień funkcji ciągłych.
			\item $C^{n}\left( R \right)$: pierścień funkcji $n$-krotnie różniczkowalnych,
			\item $C^{\infty}$ : pierścień funkcji nieskończenie wiele różniczkowalnych (gładkich).
		\end{itemize}
    \item \textbf{Pierścień Endomorfizmów} $\text{End}(A)$: (TODO JAKIS PRZYKŁAD DODAĆ)
	
	    Niech $A$ będzie grupą abelową. Wtedy
		\[
		\text{End}\left( A \right) := \{f: A \to A: f \text{ jest endomorfizmem grupy }\left( A,+ \right) \}
		.\] 
		Dla $f,g \in \text{End}\left( A \right) $ definiujemy:
		\begin{enumerate}[label=\roman*)]
			\item (Dodawanie po punktach) $f+g: A \to A$ wzorem
				\[
					\left( f+g \right) \left( a \right) := f\left( a \right) +_{\scriptscriptstyle A} g\left( a \right) 
				.\]
		        \item (Mnożenie jak składanie) $f\cdot g: A \to A$ wzorem
				\[
					\left( f\cdot g \right)\left( a \right) := f\left( a \right) \cdot_{\scriptscriptstyle A} g\left( a \right) 
				.\] 
		\end{enumerate}
		Wtedy $\left( \text{End}\left( A \right), +, \cdot  \right) $ jest pierścieniem z $\mathbf{1}=\id_{A}$, zwykle nieprzemienny.
		\begin{remark}[Po co abelowość $A$?]
		 Aby suma $f+g$ była endomorfizmem, abelowość grupy $(A, +)$ jest konieczna, by zachodził warunek $(f+g)(a+b) = (f+g)(a) + (f+g)(b)$.

		 Dlaczego?

		Dla $f,g \in \text{End}\left( A \right) $ oraz $a,b\in A$:
		\begin{align*}
			\left( f+g \right) \left( a+b \right) &\overset{\Def}{=} f\left( a+b \right) + g\left( a+b \right)\\
			\text{($f,g$ są homomorfizmami) } &= \left( f\left( a \right) + f\left( b \right)  \right) + \left( g\left( a \right) + g\left( b \right)  \right) \\
						 &= f\left( a \right) + f\left( b \right) + g\left( a \right) + g\left( b \right) 
		.\end{align*}

		Aby $f+g$ było homomorfizmem musi zachodzić równość
		\[
		f\left( b \right) + g\left( a \right) = g\left( a \right) + f\left( b \right) 
		.\] 
		Założenie o abelowośći nam to gwarantuje i dzięki temu ,,domyka'' nam się warunek bycia homomorfizmem dla $f+g$: 
		\begin{align*}
			f\left( a \right) + f\left( b \right) + g\left( a \right) + g\left( b \right) &= f\left( a \right) + g\left( a \right) + f\left( b \right) + g\left( b \right) \\
						      &\overset{(\Def +)}{=} \left( f+g \right) \left( a \right) + \left( f + g \right) \left( b \right) 
		.\end{align*}
		\end{remark} 
\end{enumerate}

\subsection*{Podstawowe własności arytmetyczne pierścieni}

\begin{fact}[Arytmetyka pierścienia]
    Niech $R$ będzie pierścieniem, $a,b \in R$, $n,m \in \mathbb{N}_{>0}$, $\alpha, \beta \in \mathbb{Z}$. Wtedy zachodzą znane prawa, m.in.:
    \begin{enumerate}
        \item Prawa potęg: $a^n a^m = a^{n+m}, \quad (a^n)^m = a^{nm}$.
	\item Własność dodawania: $\alpha a + \beta a = \left( \alpha +\beta  \right) a$.
	\item Własność mnożenia: $\alpha \left( \beta a \right) = \left( \alpha \beta  \right) a$. 
        \item Mnożenie przez zero: $a \cdot \mathbf{0} = \mathbf{0} = \mathbf{0} \cdot a$.
        \item Reguła znaków: $(-a)b = -(ab) = a(-b)$ oraz $(-a)(-b) = ab$.
        \item Rozdzielność odejmowania: $a(b-c) = ab - ac$ oraz $\left( b-c \right)a = ba - ca$
        \item Jeśli $R$ ma jedynkę, to $(-\mathbf{1})a = -a$.
        \item \textbf{Wzór skróconego mnożenia:} Jeśli $R$ jest \textbf{przemienny}, to zachodzi wzór Newtona:
        \[ (a+b)^n = \sum_{k=0}^n \binom{n}{k} a^k b^{n-k} \]
        \textit{Uwaga: Bez założenia przemienności wzór ten nie działa (np. dla macierzy $(A+B)^2 = A^2 + AB + BA + B^2 \neq A^2 + 2AB + B^2$).}
    \end{enumerate}
\end{fact}
\begin{proof}
	TODO
\end{proof} 

\begin{remark}
    Czwarta własność ($a \cdot \mathbf{0} = \mathbf{0}$) implikuje, że jeśli w pierścieniu $R$ mamy $\mathbf{0} = \mathbf{1}$, to pierścień ten składa się tylko z jednego elementu $R=\{\mathbf{0}\}$. Nazywamy go pierścieniem trywialnym lub zerowym.
\end{remark}
\begin{remark}
	Własność $A2$ (rozdzielność dodawania względem mnożenia) mówi dokładnie, że dla  $\left(\forall r\in R \right) $ funkcja
	\[
	\lambda_{r}: R \to R\text{, dana wzorem } \lambda_{r}\left( x \right) := r\cdot x
	.\] 
	jest \textbf{endomorfizmem} grupy $\left( R,+ \right) $. Bo
	\[
	\lambda_{r}\left( a + b \right) = r\cdot \left( a+b \right) = r\cdot a + r \cdot b = \lambda_{r}\left( a \right) + \lambda_{r}\left( b \right) 
	.\] 
\end{remark} 
\section{Elementy odwracalne ($R^*$)}

\begin{definition}[Element odwracalny]
    Niech $R$ będzie pierścieniem z jedynką. Element $a \in R$ nazywamy \textbf{odwracalnym} (lub jednością arytmetyczną), jeśli posiada element odwrotny względem mnożenia:
    \[ (\exists b \in R) (ab = \mathbf{1} \land ba = \mathbf{1}) \]
    Taki element $b$ oznaczamy $a^{-1}$ i jest on wyznaczony jednoznacznie.
    
    Zbiór wszystkich elementów odwracalnych oznaczamy symbolem $R^*$ (lub $U(R)$):
    \[ R^* := \{ a \in R \mid a \text{ jest odwracalny} \} \]
\end{definition}
\begin{remark}
	Niech $R$ będzie pierścieniem z $\mathbf{1}$.
	Wtedy
	\[
	\left(\forall a \in R^{\ast} \right)\left(\exists !b\in R^{\ast} \right)\left( ab = ba = \mathbf{1} \right)   
	.\] 
	Oznaczamy wtedy $b = a^{-1}$.

	\noindent \textbf{Ogólniej:}

	$\left(\forall n\in \mathbb{Z}  \right)\left(\forall a \in R^{\ast} \right)  $ mamy zdefiniowane $a^{n}$. Wtedy:
	\begin{itemize}
		\item $a^{-n} = \left( a^{-1} \right)^{n}$ jest elementem odwrotnym do $a^{n}$ w sensie mnożenia,
		\item $-na$ jest elementem odwrotnym do $na$ w sensie dodawania (czyli elementem \textbf{przeciwnym} do $na$ ).
	\end{itemize}
\end{remark}

\begin{fact}
    Para $(R^*, \cdot)$ tworzy \textbf{grupę} (zwaną grupą multiplikatywną pierścienia).
\end{fact}
\begin{proof}
	Mamy zbiór $R^{\ast}$, czyli zbiór elementów odwracalnych pierścienia $R$, oraz działanie $\cdot : R \times R \to R$ obcięte do zbioru $R^{\ast} \times R^{\ast}$, tzn $\cdot \restriction_{\scriptscriptstyle R^{\ast} \times R^{\ast} } $. W dalszej części będziemy działanie  $\cdot \restriction_{\scriptscriptstyle R^{\ast} \times R^{\ast}}$ oznaczać jako $\ast $.
	Zatem do pokazania jest:
	\begin{enumerate}[label=\arabic*)]
		\item Działanie $\ast$ jest wewnętrzne (czyli $R^{\ast} $ jest zamknięte na $\ast$).
		\item $R^{\ast}$ zawiera element neutralny działania $\ast$. 
		\item Działanie $\ast$ jest łączne.
		 \item Każdy element z $R^{\ast}$ posiada element odwrotny względem działania $\ast$
	\end{enumerate}
\end{proof} 
\begin{remark}
	$R \neq \{\mathbf{0}\} \implies \mathbf{0}\not\in R^{\ast}$
\end{remark} 
\begin{remark}
	Dodawanie $\left( + \right) $ nie jest działaniem na $R^{\ast} $, chyba, że $R = \{\mathbf{0}\} $.
\end{remark} 
\begin{definition}
	Niech $R$ będzie pierścieniem z $\mathbf{1}$. Wtedy
	\begin{enumerate}[label=\arabic*)]
		\item $\left( R,+ \right) $ jest \textbf{grupą addytywną} $R$,
		\item $\left( R^{\ast}, \cdot  \right) $ jest \textbf{grupą multiplikatywną} $R$.
	\end{enumerate}
\end{definition} 

\subsection*{Katalog przykładów grup odwracalnych $R^{\ast}$}

\begin{enumerate}[label=\arabic*)]
    % Liczby całkowite
    \item $\mathbb{Z}^{\ast} = \{-1, 1\}$.
    \begin{itemize}
        \item \textbf{Wyjaśnienie:} Element $a$ jest odwracalny, jeśli istnieje $b \in \mathbb{Z}$ takie, że $ab = 1$. W zbiorze liczb całkowitych jedynymi rozwiązaniami tego równania są pary $(1,1)$ oraz $(-1,-1)$. Dowolna inna liczba (np. $2$) wymagałaby ułamka ($1/2$), który nie należy do $\mathbb{Z}$.
    \end{itemize}

    % Reszty modulo n
    \item $(\mathbb{Z}_{n})^{\ast} = \{a \in \mathbb{Z}_{n} : \text{NWD}(a,n) = 1\} \cong \text{Aut}(\mathbb{Z}_{n})$.
    \begin{itemize}
        \item \textbf{Wyjaśnienie:} Z rozszerzonego algorytmu Euklidesa wiemy, że $\text{NWD}(a,n)=1$ jest równoważne istnieniu liczb $x, y$ takich, że $ax + ny = 1$. Przechodząc na reszty modulo $n$, otrzymujemy $ax \equiv 1 \pmod{n}$, co oznacza, że $x$ jest odwrotnością $a$. Jeśli $\text{NWD}(a,n) > 1$, to $a$ jest dzielnikiem zera i nie może być odwracalny.
    \end{itemize}

    % Ciała
    \item Dla ciał $\mathbb{Q}, \mathbb{R}, \mathbb{C}$ grupa elementów odwracalnych to po prostu zbiór elementów niezerowych:
    \[ \mathbb{Q}^{\ast} = \mathbb{Q}\setminus \{0\}, \quad \mathbb{R}^{\ast} = \mathbb{R} \setminus \{0\}, \quad \mathbb{C}^{\ast} = \mathbb{C} \setminus \{0\}. \]
    \begin{itemize}
        \item \textbf{Wyjaśnienie:} Z definicji ciała, każdy niezerowy element musi posiadać odwrotność multiplikatywną. W tych pierścieniach "dzielenie" przez dowolną liczbę różną od zera jest zawsze wykonalne i daje wynik w tym samym zbiorze.
    \end{itemize}

        % Macierze - sekcja rozbudowana i uporządkowana
    \item \textbf{Grupy macierzowe:} Niech $R$ będzie pierścieniem przemiennym z $\mathbf{1}$. 
    Definiujemy \textit{Ogólną Grupę Liniową} (\textbf{G}eneral \textbf{L}inear) :
    \[ 
        GL_{n}(R) := M_n(R)^{\ast} = \{A\in M_{n}(R) : \det(A) \in R^{\ast} \}.
    \]
    Konkretne przykłady:
    \begin{itemize}
        \item $GL_{n}(\mathbb{R}) = \{A\in M_{n}(\mathbb{R}) : \det(A) \neq 0\}$,
        \item $GL_{n}(\mathbb{Z}) = \{A\in M_{n}(\mathbb{Z}) : \det(A) \in \{-1, 1\}\}$,
        \item $GL_{n}(\mathbb{Z}_n) = \{A\in M_{n}(\mathbb{Z}_n) : \det(A) \in \mathbb{Z}_n^{\ast}\}$.
    \end{itemize}
    
    Ponadto wyróżniamy podgrupę \textit{Specjalną Grupę Liniową}:
    \[
        SL_{n}(R) = \{A \in GL_{n}(R) : \det(A) = \mathbf{1}\}.
    \]
    Zauważmy, że $SL_{n}(R) \unlhd GL_{n}(R)$. Grupa ilorazowa przez centrum to \textit{Projective Special Linear Group}:
    \[
        PSL_{n}(R) := SL_{n}(R) \big/ Z(SL_{n}(R)).
    \]
% Macierze konkretny przykład
    \item $M_{n}\left( \mathbb{R}  \right) ^{\ast} = GL_{n}\left( \mathbb{R}  \right) $. Podobnie dla $\mathbb{Q} \text{ oraz }\mathbb{C} $.
    \begin{itemize}
        \item \textbf{Wyjaśnienie:} Wzór na macierz odwrotną to $A^{-1} = \frac{1}{\det(A)} \text{adj}(A)$. Aby macierz $A^{-1}$ miała wyrazy w pierścieniu $R$, musimy móc ,,podzielić'' przez wyznacznik, co oznacza, że $\det(A)$ musi być elementem odwracalnym w $R$. 
        \item \textit{Przykład:} W $M_n(\mathbb{Z})$ macierz jest odwracalna tylko gdy $\det(A) = \pm 1$.
    \end{itemize}

    % Funkcje
    \item Niech $X$ będzie zbiorem, a $R$ pierścieniem z $\mathbf{1}$. Wtedy:
    \[ \text{Fun}(X,R)^{\ast} = \text{Fun}(X, R^{\ast}). \]
     \begin{itemize}
        \item \textbf{Wyjaśnienie:} Działania w pierścieniu funkcji są definiowane punktowo: $(f \cdot g)(x) = f(x)g(x)$. Aby funkcja $g$ była odwrotnością $f$, w każdym punkcie $x$ musi zachodzić $f(x)g(x) = \mathbf{1}_R$. To wymusza, aby dla każdego argumentu $x$, wartość $f(x)$ była elementem odwracalnym w pierścieniu $R$.
    \end{itemize}

    % Endomorfizmy
    \item Niech $A$ będzie grupą abelową. Wtedy:
    \[ \text{End}(A)^{\ast} = \text{Aut}(A). \]
    \begin{itemize}
        \item \textbf{Wyjaśnienie:} W pierścieniu endomorfizmów mnożenie to składanie odwzorowań ($f \circ g$). Elementem odwracalnym jest więc endomorfizm, który posiada odwrotność względem składania, czyli taki, który jest jednocześnie bijekcją. Endomorfizm będący bijekcją to z definicji automorfizm grupy $A$.
    \end{itemize}
\end{enumerate}

\section{Wstęp do teorii podzielności.}
 
\begin{context}
    W tym rozdziale zakładamy, że $R$ jest \textbf{pierścieniem przemiennym}.
\end{context}

\begin{definition}[Podzielność i Stowarzyszenie]
    Niech $x, y \in R$. 
    \begin{enumerate}[label=\arabic*)]
        \item Mówimy, że $x$ \textbf{dzieli} $y$ (ozn. $x \mid y$), gdy istnieje taki element $r \in R$, że $y = rx$.
        \item Mówimy, że $x$ jest \textbf{stowarzyszony} z $y$ (ozn. $x \sim y$), gdy $x \mid y$ oraz $y \mid x$.
    \end{enumerate}
\end{definition}

\begin{eg}
    W pierścieniu liczb całkowitych $\mathbb{Z}$:
    \[
     x \sim y \iff x = \pm y
    .\] 
\end{eg}
\begin{proof}
\noindent $\left( \implies \right) $ 

Ustalmy dowolny $x,y \in \mathbb{Z}$ i załóżmy, że $x\sim y$. Oznacza to, że $x \mid y$ oraz $y\mid x$, w szczególności $y = ax$ oraz  $x = by$, dla pewnych  $a,b \in \mathbb{Z}$. Przekształcając równoważnie dane równości, otrzymujemy, że $y = \left( ab \right) x$ oraz $x = \left( ba \right) y$. Łącząc obie równości, mamy $y = \left( abba \right)y $ oraz $x = \left( baab \right) x$. Ponieważ $\mathbb{Z}$ jest przemienny więc $abba = baab = abab = \left( ab \right) ^2$. Zatem $x = \left( ab \right) ^2x$ oraz $y = \left( ab \right) ^2y$, czyli $\left( ab \right) ^2 = 1$, co w $\mathbb{Z}$ oznacza, że $ab = \pm 1$. Zatem $y = \pm x$ oraz  $x = \pm y$. $\diamond $

\noindent $\left( \impliedby \right) $ 

Ustalmy dowolne $x,y \in \mathbb{Z}$ i załóżmy, że $x = \pm y$. Oznacza to, że  $x = y$ lub  $x = -y$.
 \begin{itemize}
	\item Przypadek $x = y$ oznacza, że  $x = 1\cdot y$ oraz $y = 1\cdot x$, czyli $x \mid y$ oraz $y \mid x$, więc na mocy definicji $\sim  $,  $x\sim y $.
	\item Przypadek $x = -y$ ($\iff -x = y$) oznacza, że $x = \left( -1 \right) y$ oraz $y = \left( -1 \right) x$, czyli $x \mid y$ oraz $y \mid x$, więc z definicji $\sim  $, $x\sim y $.
\end{itemize}
\end{proof} 
\begin{fact}
    Niech $x, y, z \in R$.
    \begin{enumerate}[label=\arabic*)]
        \item Jeśli $R$ ma jedynkę, relacja $\sim$ jest relacją równoważności.
        \item $(\exists n \in R^{\ast})(y = nx) \implies y \sim x$.
    \end{enumerate}
\end{fact}
\begin{proof}
	
	Ustalmy dowolne $x,y,z \in R$.

	\noindent $\bm{ 1) } $ 

	Załóżmy, że $R$ jest pierścienie z jedynką. Pokażemy, że relacja  $\sim$ jest relacją równoważności.
	\begin{itemize}
		\item Zwrotność: Ponieważ $x = 1\cdot x$ (istnieje jedynka), więc $x \mid x$, zatem $x \sim x$.
		\item Symetryczność: Załóżmy, że $x \sim y $, wtedy $x \mid y$ oraz $y \mid x$, co jest równoważne $y \mid x$ oraz $x \mid y$, zatem $y \sim x$.
		\item Przechodniość: Załóżmy, że $x\sim y $ oraz $y \sim z$. Oznacza to, że 
			\[
			x \mid y \land y \mid x \land y \mid z \land z \mid y
			.\] 
		W szczególności
		\begin{itemize}
			\item $x = ay$,
			\item $y = bx$,
			\item $y = cz$,
			\item $z = dy$,
		\end{itemize}
		dla pewnych $a,b,c,d \in R$. \\
		Ponieważ $x = ay$ i  $y = cz$ więc  $x = \left( ac \right) z$, czyli $x \mid z$.\\
		Analogicznie $y = bx$ i  $z = dy$, co nam daje  $z = \left( db \right) x$, czyli $z \mid x$.

		Zatem $x \mid z$ oraz $z \mid x$, co z definicji oznacza, że $x\sim z $. $\diamond $
	\end{itemize}

	\noindent $\bm{ 2) } $

	TODO

	Załóżmy, że istnieje odwracalne $n$ takie, że $y = nx$, tzn  $\left(\exists n \in R^{\ast}  \right) \left(  \right)  $


\end{proof} 
\subsection*{Dzielnik zera oraz dziedzina całkowitości.}

W pierścieniach może zdarzyć się patologia: iloczyn dwóch niezerowych liczb daje zero.

\begin{definition}[Dzielnik zera]
    Element $x \in R \setminus \{\mathbf{0}\}$ nazywamy \textbf{dzielnikiem zera}, jeśli istnieje taki $y \in R \setminus \{\mathbf{0}\}$, że $xy = \mathbf{0}$ (lub $yx = \mathbf{0} $ ).
\end{definition}

\begin{eg}
	\begin{itemize}
		\item W pierścieniu $\mathbb{Z}_4$: \[2 \cdot 2 = 4 \equiv 0 \pmod 4.\] Zatem $2$ jest dzielnikiem zera.
		\item W pierścieniu $M_2(\mathbb{R})$ (macierz nilpotentna): \[\begin{pmatrix} 0 & 1 \\ 0 & 0 \end{pmatrix} \begin{pmatrix} 0 & 1 \\ 0 & 0 \end{pmatrix} = \begin{pmatrix} 0 & 0 \\ 0 & 0 \end{pmatrix}.\]
	\end{itemize}
\end{eg}

\begin{definition}[Dziedzina Całkowitości]
	Pierścień $R$ nazywamy \textbf{dziedziną całkowitości} (ang. \textit{integral domain}) (w skrócie dziedziną), gdy spełnia trzy warunki:
    \begin{enumerate}
        \item Jest przemienny z jedynką.
        \item Jest nietrywialny ($\mathbf{1} \neq \mathbf{0}$).
        \item \textbf{Nie posiada dzielników zera} (tzn. $ab=0 \implies a=0 \lor b=0$).
    \end{enumerate}
\end{definition}
\begin{intuition}
    Dziedziny to pierścienie, w których arytmetyka przypomina tę w liczbach całkowitych $\mathbb{Z}$. Jedną z ważniejszych własności dziedzin jest \textbf{prawo skracania}:
    \[ \text{Jeśli } ax = ay \text{ i } a \neq \mathbf{0}, \text{ to } x = y. \]

    Tak naprawdę, jest to warunek równoważny na bycie dziedziną, o czym mówi poniższe twierdzenie.
\end{intuition}
\begin{theorem}	
	Pierścień $R$ przemienny z jedynką jest dziedziną całkowitości wtedy i tylko wtedy, kiedy zachowane jest praw skracania, tj
	\[
	\left(\forall a,b,c \in R \right)\left( \left(ab = ac \text{ i } a \neq 0 \right) \implies b = c\right)  
	.\] 
\end{theorem} 
\begin{proof}
	TODO
\end{proof} 
\begin{fact}
    Niech $R$ będzie pierścieniem przemiennym z $\mathbf{1}$ oraz $x,y \in R$. Wtedy:
    \begin{enumerate}[label=\arabic*)]
        \item $R$ jest dziedziną $\implies \left( x\sim y \iff \exists n \in R^{\ast} \text{ taki, że } y = nx \right)$.
        \item $x\in R^{\ast} \implies x$ nie jest dzielnikiem zera.
    \end{enumerate}
\end{fact}
\begin{proof}
    1) ($\Leftarrow$) Załóżmy, że $y = nx$ dla pewnego $n \in R^{\ast}$. Z definicji podzielności $x \mid y$. Ponieważ $n$ jest odwracalny, istnieje $n^{-1} \in R$, zatem $x = n^{-1}y$, co oznacza $y \mid x$. Stąd z definicji $x \sim y$. \\
    ($\Rightarrow$) Załóżmy, że $R$ jest dziedziną i $x \sim y$. Jeśli $x = \mathbf{0}$, to z $x \mid y$ wynika $y = \mathbf{0}$, więc teza zachodzi dla $n=\mathbf{1}$. Jeśli $x \neq \mathbf{0}$, to z relacji stowarzyszenia mamy $y = nx$ oraz $x = my$ dla pewnych $n, m \in R$[cite: 48]. Podstawiając, otrzymujemy $x = m(nx) = (mn)x$, co można zapisać jako $(mn - \mathbf{1})x = \mathbf{0}$. Ponieważ $R$ jest dziedziną i $x \neq \mathbf{0}$, z braku dzielników zera wynika, że $mn - \mathbf{1} = \mathbf{0}$, czyli $mn = \mathbf{1}$. Zatem $n \in R^{\ast}$.

    2) Niech $x \in R^{\ast}$. Przypuśćmy nie wprost, że $x$ jest dzielnikiem zera. Wtedy $x \neq \mathbf{0}$ i istnieje $y \neq \mathbf{0}$ takie, że $xy = \mathbf{0}$. Mnożąc obie strony równania przez $x^{-1}$, otrzymujemy $x^{-1}(xy) = x^{-1} \cdot \mathbf{0}$, co upraszcza się do $\mathbf{1} \cdot y = \mathbf{0}$, czyli $y = \mathbf{0}$[cite: 100]. Otrzymana sprzeczność z założeniem $y \neq \mathbf{0}$ dowodzi, że $x$ nie może być dzielnikiem zera.
\end{proof}

\begin{remark}
        Zauważmy, że istnieją pierścienie z $\mathbf{1}$, które nie są dziedzinami i w których relacja $x \sim y$ nie musi oznaczać istnienia elementu odwracalnego $n$ łączącego te elementy (przykładem może być pierścień funkcji ciągłych $C[0,1]$). 

	Temat ten zostanie poruszony na ostatnim wykładzie (14).
\end{remark}

\section{Wprowadzenie ciał.}

\begin{definition}[Pierścień z dzieleniem i Ciało]
    Niech $R$ będzie pierścieniem z jednością $\mathbf{1}$.
    \begin{enumerate}[label=\arabic*)]
        \item $R$ nazywamy \textbf{pierścieniem z dzieleniem} (lub pierścieniem z dzielnikiem), gdy każdy jego niezerowy element jest odwracalny, czyli:
        \[ R^{\ast} = R \setminus \{\mathbf{0}\} \]
        Warunek ten implikuje, że pierścień jest nietrywialny ($\mathbf{0} \neq \mathbf{1}$).
        \item $R$ nazywamy \textbf{ciałem}, gdy jest przemiennym pierścieniem z dzieleniem.
    \end{enumerate}
\end{definition}
\begin{intuition}	
    Pierścień przemienny z jedynką $R$ ($\mathbf{1} \neq \mathbf{0}$) jest \textbf{ciałem}, gdy każdy niezerowy element jest odwracalny.
    \[ R^* = R \setminus \{ \mathbf{0} \} \]
\end{intuition} 

\noindent \textbf{Przykłady i klasyfikacja struktur:}
\begin{enumerate}[label=\arabic*)]
    \item $\mathbb{Q}, \mathbb{R}, \mathbb{C}$ są ciałami.
    \item $\mathbb{Z}$ jest dziedziną, ale \textbf{nie} jest ciałem, ponieważ dla $a \in \mathbb{Z} \setminus \{-1, 1\}$ nie istnieje element odwrotny wewnątrz $\mathbb{Z}$.
    
    \item \textbf{Kwaterniony ($\mathbb{H}$):} Jest to pierścień z dzieleniem (skew field), który \textbf{nie} jest ciałem. 
    \begin{itemize}
        \item Każdy niezerowy kwaternion $q = a + bi + cj + dk$ posiada element odwrotny $q^{-1} = \frac{\bar{q}}{|q|^2}$, gdzie $\bar{q} = a - bi - cj - dk$.
        \item Mnożenie nie jest przemienne: zachodzą relacje $ij = k$ oraz $ji = -k$, co wyklucza $\mathbb{H}$ z definicji ciała.
    \end{itemize}

    \item \noindent \textbf{Oznaczenie: } $\mathbb{F}_{p}: \mathbb{Z}_{p}$ ciało p-elementowe. 

	\textbf{Wyjaśnienie oznaczenia $\mathbb{F}_p$:} 
	W literaturze symbol $\mathbb{F}_p$ (od ang. \textit{field}) oznacza abstrakcyjne ciało o charakterystyce $p$ mające $p$ elementów[cite: 60]. W Twoich notatkach jest ono utożsamiane z $\mathbb{Z}_p$, ponieważ zbiór reszt modulo $p$ (gdzie $p$ jest liczbą pierwszą) dostarcza nam konkretnej, konstruktywnej realizacji takiego ciała.

    \item Niech $n \in \mathbb{N}_{>0}$. Następujące warunki są równoważne:
    \begin{enumerate}[label=\roman*)]
        \item $\mathbb{Z}_n$ jest ciałem.
        \item $\mathbb{Z}_n$ jest dziedziną.
        \item $n$ jest liczbą pierwszą.
    \end{enumerate}
    \textit{Ważne:} Równoważność $i) \iff ii)$ zachodzi dla wszystkich pierścieni przemiennych \textbf{skończonych}:
\end{enumerate}
\begin{theorem}
Każda skończona dziedzina całkowitości jest ciałem.
\end{theorem}

\begin{proof}[Szkic dowodu]
Niech $R = \{x_1, x_2, \dots, x_n\}$ będzie skończoną dziedziną całkowitości. 
Aby wykazać, że $R$ jest ciałem, musimy udowodnić, że każdy element $a \in R \setminus \{0\}$ posiada odwrotność multiplikatywną.

\begin{enumerate}
    \item \textbf{Konstrukcja odwzorowania:} Dla ustalonego $a \neq 0$ rozważmy odwzorowanie mnożenia lewostronnego $\lambda_a: R \to R$ dane wzorem:
    \[ \lambda_a(x) = ax \]
    
    \item \textbf{Iniekcja (Różnowartościowość):} Załóżmy, że $\lambda_a(x) = \lambda_a(y)$. Wtedy $ax = ay$, co implikuje $a(x-y) = 0$. 
    Ponieważ $R$ jest dziedziną i $a \neq 0$, nie istnieją dzielniki zera, więc musi zachodzić $x-y = 0$, czyli $x=y$. 
    Odwzorowanie $\lambda_a$ jest zatem iniekcją.
    
    \item \textbf{Surjekcja i Jedynka:} Ponieważ $R$ jest zbiorem skończonym, każda iniekcja ze zbioru w ten sam zbiór jest automatycznie surjekcją (zasada szufladkowa Dirichleta). 
    Skoro $\lambda_a$ jest surjekcją, to w obrazie tego odwzorowania musi znajdować się element neutralny mnożenia $\mathbf{1} \in R$.
    
    \item \textbf{Konkluzja:} Istnieje zatem taki element $b \in R$, że $\lambda_a(b) = \mathbf{1}$, co oznacza $ab = \mathbf{1}$. 
    Z przemienności pierścienia wynika, że $b$ jest szukaną odwrotnością elementu $a$.
\end{enumerate}
Zatem każdy niezerowy element $a$ jest odwracalny, co czyni $R$ ciałem.
\end{proof}

\begin{remark}[Hierarchia struktur]
Mamy następujący ciąg inkluzji (każda struktura po prawej jest ogólniejsza):
\[ \text{Ciało} \implies \text{Dziedzina} \implies \text{Pierścień przemienny z } \mathbf{1} \implies \text{Pierścień z } \mathbf{1} \]
Ale
	\[
	\mathbb{R} \nLeftarrow \mathbb{Z} \nLeftarrow \mathbb{Z} _{n} \nLeftarrow M_{n}\left( \mathbb{R}  \right) 
	.\] 
Implikacje w drugą stronę (z prawa na lewo) zazwyczaj nie zachodzą, co pokazują poniższe kontrprzykłady:
\begin{itemize}
	\item \textbf{Pierścień macierzy $M_n(\mathbb{R})$} dla $n > 1$ posiada $\mathbf{1}$, ale nie jest przemienny.
	\item \textbf{Pierścień $\mathbb{Z}_n$} dla $n$ złożonego (np. $\mathbb{Z}_4$) jest przemienny, ale nie jest dziedziną, bo posiada dzielniki zera ($2 \cdot 2 \equiv 0 \pmod 4$).
	\item \textbf{Liczby całkowite $\mathbb{Z}$} są dziedziną, ale nie są ciałem, bo tylko $1$ i $-1$ są odwracalne.
\end{itemize}
\end{remark}

\begin{definition}[Podciało]
Niech $L$ będzie ciałem. Podzbiór $K \subseteq L$ nazywamy \textbf{podciałem} ciała $L$ (ozn. $K \underset{\text{podc}}{\subseteq} L$), gdy spełnione są następujące warunki:
\begin{enumerate}
	\item $K$ jest podpierścieniem ciała $L$:  \[ K\underset{\scriptscriptstyle\text{podp}}{\subseteq} L .\]
    \item $\mathbf{1}_L \in K$ (jedynka ciała $L$ należy do podzbioru $K$).
    \item Dla każdego niezerowego elementu $a \in K \setminus \{0\}$, jego element odwrotny względem mnożenia również należy do $K$ ($a^{-1} \in K$):
	    \[
	     \left(\forall a \in K\setminus \{\mathbf{0}\}  \right) \left( a^{-1}\in K \right)
	    .\] 
\end{enumerate}
\end{definition}

\begin{remark}
Jeśli $K$ jest podciałem $L$, to struktura $\left( K, +\restriction_{\scriptscriptstyle K\times K}, \cdot\restriction_{\scriptscriptstyle K\times K} \right) $  z działaniami zawężonymi do zbioru $K$ jest ciałem w sensie ogólnej definicji.
\end{remark}

\noindent \textbf{Przykłady:}
\begin{enumerate}[label=\arabic*)]
    	\item \textbf{Ciąg inkluzji:} $\mathbb{Q} \underset{\text{podc}}{\subseteq} \mathbb{R} \underset{\text{podc}}{\subseteq} \mathbb{C}$. Każdy z tych zbiorów jest podciałem kolejnego, ponieważ zachowuje operacje dodawania i mnożenia oraz zawiera odwrotności wszystkich swoich niezerowych elementów.
    	\item \textbf{Przypadek pierścienia liczb całkowitych:} $\mathbb{Z}$ jest podpierścieniem ciała $\mathbb{Q}$ ($\mathbb{Z} \underset{\text{podp}}{\subseteq} \mathbb{Q}$), ale \textbf{nie jest podciałem} $\mathbb{Q}$. 

    \textit{Uzasadnienie:} Choć $\mathbb{Z}$ zawiera jedynkę, to nie zawiera odwrotności swoich elementów (poza $1$ i $-1$). Przykładowo $2 \in \mathbb{Z}$, ale jego odwrotność multiplikatywna $1/2$ nie należy do zbioru liczb całkowitych. Istotnie, $\mathbb{Z} $ jest ,,zbyt małe'' aby być ciałem, pomimo iż świetnie radzi sobie jako pierścień.
\end{enumerate}

